% GENERATED FILE. Edit canonical lessons, then run npm run generate:books.
% canonical-source-hash: fnv1a64:4dd54a4d5ff547b4
% canonical-lessons: RU-C32-pain, RU-C32-illness, RU-C32-medicine, RU-C32-doctor, RU-C32-hospital

\chapter{Pain, Illness, Medicine, Doctor}
\label{ch:ru-pain-illness-medicine-doctor}
\hlchaptermodality{32}

\begin{chapteropening}
\textbf{By the end of this chapter:} \emph{I can name pain, an illness, medicine, a doctor and a hospital.}

Complete the last lesson of chapter 32: I can name pain, an illness, medicine, a doctor and a hospital.
\end{chapteropening}

\section[\texorpdfstring{\ru{боль} (bol)}{боль (bol)}]{\ru{боль} (bol) — pain}

\label{lesson:RU-C32-pain}



\begin{quote}
Before the new one: say the Russian for the face, then the Russian for blood.
\end{quote}

\subsection*{You'll want to know: \ru{боль}}
\textbf{\ru{боль}} (\emph{bol}) — \textquotedblleft{}pain\textquotedblright{}.

\begin{cousinweb}
\textbf{\ru{боль}} $\leftarrow$ Proto-Slavic \emph{*bolĭ} — inherited, with the verb \emph{\ru{болеть}}, \textquotedblleft{}to ache, to be ill\textquotedblright{}. Vasmer sets it beside Gothic \emph{balwjan}, \textquotedblleft{}to torment\textquotedblright{}, and Old English \emph{bealu}, \textquotedblleft{}evil\textquotedblright{}, the \emph{bale} of \emph{baleful}, but notes that Berneker rejected the link.
\end{cousinweb}

\begin{grammarlens}[title={What you've built}]
One more word for this chapter.
\end{grammarlens}

\subsection*{Your turn}
\begin{itemize}
\raggedright
  \item \emph{Say it:} \emph{bol}
  \item \emph{Say it:} \emph{bol}, once more
  \item \emph{Say it:} say \emph{bol}
  \item \emph{Recall:} say the Russian for the face, then the Russian for blood, then say \emph{bol} again
\end{itemize}

\begin{tcolorbox}[breakable,colback=teal!4,colframe=teal!35!black,title={Before you move on}]
What does \ru{боль} mean? (\textquotedblleft{}Pain\textquotedblright{}.) Say it once more.
\end{tcolorbox}
\section[\texorpdfstring{\ru{болезнь} (bolézn)}{болезнь (bolézn)}]{\ru{болезнь} (bolézn) — an illness}

\label{lesson:RU-C32-illness}



\begin{quote}
Before the new one: say the Russian for blood, then the Russian for pain.
\end{quote}

\subsection*{You'll want to know: \ru{болезнь}}
\textbf{\ru{болезнь}} (\emph{bolézn}) — \textquotedblleft{}an illness\textquotedblright{}. Built on \textbf{\ru{боль}}.

\begin{cousinweb}
\textbf{\ru{болезнь}} — built from \emph{\ru{болеть}}, \textquotedblleft{}to be ill\textquotedblright{}, the verb of \emph{\ru{боль}}, \textquotedblleft{}pain\textquotedblright{}, with the old noun suffix \emph{-\ru{знь}} that also makes \emph{\ru{жизнь}}, \textquotedblleft{}life\textquotedblright{}.
\end{cousinweb}

\begin{grammarlens}[title={What you've built}]
One more word for this chapter.
\end{grammarlens}

\subsection*{Your turn}
\begin{itemize}
\raggedright
  \item \emph{Say it:} \emph{bolézn}
  \item \emph{Say it:} \emph{bolézn}, once more
  \item \emph{Say it:} say \emph{bolézn}
  \item \emph{Recall:} say the Russian for blood, then the Russian for pain, then say \emph{bolézn} again
\end{itemize}

\begin{tcolorbox}[breakable,colback=teal!4,colframe=teal!35!black,title={Before you move on}]
What does \ru{болезнь} mean? (\textquotedblleft{}An illness\textquotedblright{}.) Say it once more.
\end{tcolorbox}
\section[\texorpdfstring{\ru{лекарство} (lekárstvo)}{лекарство (lekárstvo)}]{\ru{лекарство} (lekárstvo) — medicine}

\label{lesson:RU-C32-medicine}



\begin{quote}
Before the new one: say the Russian for pain, then the Russian for an illness.
\end{quote}

\subsection*{You'll want to know: \ru{лекарство}}
\textbf{\ru{лекарство}} (\emph{lekárstvo}) — \textquotedblleft{}medicine\textquotedblright{}.

\begin{cousinweb}
\textbf{\ru{лекарство}} is \emph{\ru{лекарь}}, an old word for \textquotedblleft{}doctor\textquotedblright{}, with the noun suffix \emph{-\ru{ство}}. The root \emph{\ru{лек}-}, \textquotedblleft{}cure\textquotedblright{}, came into Slavic from Germanic; Vasmer calls the loan certain and compares Gothic \emph{lēkeis}, \textquotedblleft{}doctor\textquotedblright{}. English \emph{leech}, an old word for \textquotedblleft{}physician\textquotedblright{}, is from that Germanic word.
\end{cousinweb}

\begin{grammarlens}[title={What you've built}]
One more word for this chapter.
\end{grammarlens}

\subsection*{Your turn}
\begin{itemize}
\raggedright
  \item \emph{Say it:} \emph{lekárstvo}
  \item \emph{Say it:} \emph{lekárstvo}, once more
  \item \emph{Say it:} say \emph{lekárstvo}
  \item \emph{Recall:} say the Russian for pain, then the Russian for an illness, then say \emph{lekárstvo} again
\end{itemize}

\begin{tcolorbox}[breakable,colback=teal!4,colframe=teal!35!black,title={Before you move on}]
What does \ru{лекарство} mean? (\textquotedblleft{}Medicine\textquotedblright{}.) Say it once more.
\end{tcolorbox}
\section[\texorpdfstring{\ru{врач} (vrach)}{врач (vrach)}]{\ru{врач} (vrach) — a doctor}

\label{lesson:RU-C32-doctor}



\begin{quote}
Before the new one: say the Russian for an illness, then the Russian for medicine.
\end{quote}

\subsection*{You'll want to know: \ru{врач}}
\textbf{\ru{врач}} (\emph{vrach}) — \textquotedblleft{}a doctor\textquotedblright{}.

\begin{cousinweb}
\textbf{\ru{врач}} — a Church Slavonic word, Old Church Slavonic \emph{vračĭ}; its \emph{-\ru{ра}-}, where East Slavic would say \emph{-\ru{оро}-}, marks it. Vasmer derives it from the root of \emph{\ru{врать}}, \textquotedblleft{}to talk nonsense\textquotedblright{}, and \emph{\ru{ворчать}}, \textquotedblleft{}to grumble\textquotedblright{}, so the first sense was \textquotedblleft{}charmer, healer by spells\textquotedblright{}; Bulgarian \emph{\ru{врач}} still means \textquotedblleft{}sorcerer\textquotedblright{}.
\end{cousinweb}

\begin{grammarlens}[title={What you've built}]
One more word for this chapter.
\end{grammarlens}

\subsection*{Your turn}
\begin{itemize}
\raggedright
  \item \emph{Say it:} \emph{vrach}
  \item \emph{Say it:} \emph{vrach}, once more
  \item \emph{Say it:} say \emph{vrach}
  \item \emph{Recall:} say the Russian for an illness, then the Russian for medicine, then say \emph{vrach} again
\end{itemize}

\begin{tcolorbox}[breakable,colback=teal!4,colframe=teal!35!black,title={Before you move on}]
What does \ru{врач} mean? (\textquotedblleft{}A doctor\textquotedblright{}.) Say it once more.
\end{tcolorbox}
\section[\texorpdfstring{\ru{больница} (bolnítsa)}{больница (bolnítsa)}]{\ru{больница} (bolnítsa) — a hospital}

\label{lesson:RU-C32-hospital}



\begin{quote}
Before the new one: say the Russian for medicine, then the Russian for a doctor.
\end{quote}

\subsection*{You'll want to know: \ru{больница}}
\textbf{\ru{больница}} (\emph{bolnítsa}) — \textquotedblleft{}a hospital\textquotedblright{}. Built on \textbf{\ru{боль}} as well.

\begin{cousinweb}
\textbf{\ru{больница}} is \emph{\ru{больной}}, \textquotedblleft{}sick; a patient\textquotedblright{}, with \emph{-\ru{ница}}, a suffix that names places: a place for the sick. The root is that of \emph{\ru{боль}}, \textquotedblleft{}pain\textquotedblright{}. Vasmer relates \emph{\ru{боль}} to Old English \emph{bealu}, \textquotedblleft{}harm\textquotedblright{}.
\end{cousinweb}

\begin{grammarlens}[title={What you've built}]
One more word for this chapter.
\end{grammarlens}

\subsection*{Your turn}
\begin{itemize}
\raggedright
  \item \emph{Say it:} \emph{bolnítsa}
  \item \emph{Say it:} \emph{bolnítsa}, once more
  \item \emph{Say it:} say \emph{bolnítsa}
  \item \emph{Recall:} say the Russian for medicine, then the Russian for a doctor, then say \emph{bolnítsa} again
\end{itemize}

\begin{tcolorbox}[breakable,colback=teal!4,colframe=teal!35!black,title={Before you move on}]
What does \ru{больница} mean? (\textquotedblleft{}A hospital\textquotedblright{}.) Say it once more.
\end{tcolorbox}
